\chapter*{Resumo}
\addcontentsline{toc}{chapter}{Resumo}

O HTTP/3 usa QUIC/UDP para reduzir latência e melhorar resiliência em redes instáveis. Node.js não oferece suporte nativo estável, então este TCC examina se é viável integrar HTTP/3 em APIs REST em Node.js, quais desafios surgem, como o desempenho se compara a outras versões e como fazer essa integração na prática. O trabalho combina levantamento bibliográfico, revisão de literatura e experimentos com uma API real. Construímos uma API para testes de carga, expondo três caminhos: HTTP/1.1 direto, e HTTP/2 e HTTP/3 através de um proxy reverso NGINX. Usamos k6 para carregar HTTP/1.1 e HTTP/2, e um cliente QUIC em Docker para HTTP/3. Rodamos dois conjuntos de testes: um com a rede degradada (usando tc netem, GET /bench/ping, 50 usuários virtuais em quatro perfis) e outro com a rede limpa (testando diferentes tamanhos de payload, funcionalidade echo e 10, 50 ou 100 usuários virtuais simultâneos). Cada teste teve dez rodadas, com resultados resumidos por mediana e intervalo interquartil. O HTTP/3 funcionou sem falhas em todos os testes (120 execuções com rede degradada, 210 com rede limpa), o que mostra que é viável usar um proxy na borda e backend em HTTP/1.1. Porém, geradores de carga diferentes, topologias assimétricas (HTTP/1.1 direto versus HTTP/2 e HTTP/3 via NGINX) e limitações do cliente QUIC dificultaram conclusões gerais sobre qual versão é realmente melhor. Os resultados descrevem este arranjo específico. A integração via proxy reverso é prática e não exigiu mudanças na aplicação, mas a negociação Alt-Svc e a alternância entre versões não foram testadas. O código está aberto e os roteiros de coleta em \url{https://github.com/mendeslian/QUICBench-Node} permitem replicar e estender os testes.

\noindent\textbf{Palavras-chave:} HTTP/3; QUIC; Node.js; APIs REST; viabilidade; desempenho.
